\SourcePage{84}{87}
\section{General Properties of One-Dimensional Motion}

When a particle's potential energy depends on just one coordinate $x$, the
wave function may be sought as a product of a function of $y,z$ and a function
of $x$ alone. The first is determined by the free-motion Schrödinger equation,
and the second by the one-dimensional Schrödinger equation
\begin{equation}
  \frac{d^2\psi}{dx^2}+\frac{2m}{\hbar^2}[E-U(x)]\psi=0.
  \tag{21.1}
\end{equation}
The problem of motion in a field with potential energy
$U(x,y,z)=U_1(x)+U_2(y)+U_3(z)$, separated into a sum of functions each
depending on only one coordinate, plainly reduces to the same kind of
one-dimensional equations. In \secsref{22}{24} we shall consider several specific
examples of such ``one-dimensional'' motion. First we establish some of its
general properties.

We begin by showing that every discrete-spectrum energy level in a
one-dimensional problem is nondegenerate. Suppose the contrary, and let
$\psi_1$ and $\psi_2$ be two different eigenfunctions belonging to the same
energy value. Since both satisfy the same equation (21.1),
\[
  \frac{\psi_1''}{\psi_1}
  =\frac{2m}{\hbar^2}(U-E)
  =\frac{\psi_2''}{\psi_2},
\]
or $\psi_2\psi_1''-\psi_1\psi_2''=0$ (a prime denotes differentiation with
respect to $x$). Integration gives
\begin{equation}
  \psi_1\psi_2'-\psi_2\psi_1'=\operatorname{const}.
  \tag{21.2}
\end{equation}
\SourcePage{85}{88}
Since $\psi_1=\psi_2=0$ at infinity, the constant must be zero, and hence
\[
  \psi_1\psi_2'-\psi_2\psi_1'=0,
\]
or $\psi_1'/\psi_1=\psi_2'/\psi_2$. A second integration gives
$\psi_1=\operatorname{const}\cdot\psi_2$; the two functions are therefore
essentially identical.

For the discrete-spectrum wave functions $\psi_n(x)$, the following result,
known as the \emph{oscillation theorem}, can be stated: the function
$\psi_n(x)$ corresponding to the $(n+1)$th eigenvalue $E_n$ in increasing
order vanishes $n$ times at finite $x$\footnote{If the particle can occupy
only a bounded interval of the $x$ axis, the zeros of $\psi_n(x)$ inside that
interval are to be counted.}.

Suppose that $U(x)$ tends to finite limits as $x\to\pm\infty$ (it need not in
any way be monotonic). Take $U(+\infty)$ as the zero of energy, so that
$U(+\infty)=0$, and denote $U(-\infty)$ by $U_0$, with $U_0>0$. The discrete
spectrum lies in the range of energies for which the particle cannot escape
to infinity. The energy must therefore be less than both limiting values and,
in particular, negative:
\begin{equation}
  E<0.
  \tag{21.3}
\end{equation}
Of course $E>U_{\min}$ must also hold in every case, so $U(x)$ must possess at
least one minimum with $U_{\min}<0$.

Now consider the positive energies below $U_0$:
\begin{equation}
  0<E<U_0.
  \tag{21.4}
\end{equation}
Here the spectrum is continuous and the particle motion in the corresponding
stationary states is infinite, with the particle escaping toward
$x=+\infty$. Every energy eigenvalue in this part of the spectrum is likewise
nondegenerate. To see this, it suffices to observe that the proof given above
for the discrete spectrum requires $\psi_1$ and $\psi_2$ to vanish at only
one of the two infinities (here they vanish as $x\to-\infty$).

At sufficiently large positive $x$, $U(x)$ can be neglected in the
Schrödinger equation (21.1):
\[
  \psi''+\frac{2m}{\hbar^2}E\psi=0.
\]
\SourcePage{86}{89}
This equation has real solutions in the form of a standing plane wave,
\begin{equation}
  \psi=a\cos(kx+\delta),
  \tag{21.5}
\end{equation}
where $a$ and $\delta$ are constants, and the ``wave vector'' is
$k=p/\hbar=\sqrt{2mE}/\hbar$. This formula gives the asymptotic form, as
$x\to+\infty$, of the wave functions for the nondegenerate energy levels in
the continuous-spectrum interval (21.4). At large negative $x$, the
Schrödinger equation becomes
\[
  \psi''-\frac{2m}{\hbar^2}(U_0-E)\psi=0.
\]
The solution which does not become infinite as $x\to-\infty$ is
\begin{equation}
  \psi=be^{\kappa x},
  \qquad
  \kappa=\frac{1}{\hbar}\sqrt{2m(U_0-E)}.
  \tag{21.6}
\end{equation}
This is the asymptotic wave function as $x\to-\infty$. The wave function thus
decays exponentially into the region in which $E<U_0$.

Finally, when
\begin{equation}
  E>U_0,
  \tag{21.7}
\end{equation}
the spectrum is continuous and the motion is infinite in both directions.
Every level in this part of the spectrum is doubly degenerate. The reason is
that the corresponding wave functions are determined by the second-order
equation (21.1), and both independent solutions meet the required conditions
at infinity (whereas, for example, in the preceding case one solution became
infinite as $x\to-\infty$ and had to be rejected). As $x\to+\infty$, the
asymptotic wave function is
\begin{equation}
  \psi=a_1e^{ikx}+a_2e^{-ikx},
  \tag{21.8}
\end{equation}
and similarly as $x\to-\infty$. The term containing $e^{ikx}$ corresponds to
a particle moving to the right, and that containing $e^{-ikx}$ to a particle
moving to the left.

Suppose $U(x)$ is even, $U(-x)=U(x)$. The Schrödinger equation (21.1) is then
unchanged on reversing the sign of the coordinate. Hence, if $\psi(x)$ is a
solution, $\psi(-x)$ is also a solution and differs from $\psi(x)$ only by a
constant factor: $\psi(-x)=c\psi(x)$. Reversing $x$ once more gives
$\psi(x)=c^2\psi(x)$, so $c=\pm1$. Thus, for a potential energy symmetric
about $x=0$, stationary-state wave functions may be either even,
$\psi(-x)=\psi(x)$,%
\SourcePage{87}{90}
or odd, $\psi(-x)=-\psi(x)$\footnote{This argument assumes that the stationary
state is nondegenerate, that is, that its motion is not infinite in both
directions. Otherwise the two wave functions belonging to the given energy
level may transform into one another when the sign of $x$ is reversed. Even
then, the stationary-state wave functions are not necessarily
even or odd, but definite parity can always be achieved by forming suitable
linear combinations of the two solutions.}. In particular, the ground-state
wave function is even: it cannot have a node, whereas an odd function
necessarily vanishes at $x=0$, since $\psi(0)=-\psi(0)=0$.

There is a simple method for normalizing one-dimensional-motion wave functions
in the continuous spectrum, which determines the normalization coefficient
directly from the asymptotic wave function at large $|x|$.

Consider the wave function for motion infinite in one direction
($x\to+\infty$). The normalization integral diverges as $x\to\infty$ (at
$x\to-\infty$ the function decays exponentially, so the integral converges
rapidly). To determine the normalization constant, one may therefore replace
$\psi$ by its asymptotic value for large positive $x$ and take any finite
value of $x$, say zero, as the lower limit of integration; this is equivalent
to discarding a finite contribution beside a divergent one. We now
show that a wave function normalized by
\begin{equation}
  \int\psi_p^*\psi_{p'}\,dx
  =\delta\left(\frac{p-p'}{2\pi\hbar}\right)
  =2\pi\hbar\delta(p-p')
  \tag{21.9}
\end{equation}
(where $p$ is the particle momentum at infinity) must have the asymptotic form
(21.5) with $a=2$:
\begin{equation}
  \psi_p\approx2\cos(kx+\delta)
  =e^{i(kx+\delta)}+e^{-i(kx+\delta)}.
  \tag{21.10}
\end{equation}
Since we are not seeking to verify mutual orthogonality of functions
corresponding to different $p$, when substituting (21.10) into the
normalization integral we take $p$ and $p'$ arbitrarily close and may put
$\delta=\delta'$ (in general, $\delta$ is a function of $p$). We then retain
in the integrand only the terms which diverge when $p=p'$; equivalently, we
drop terms containing factors $e^{\pm i(k+k')x}$. We thus obtain
\[
  \int\psi_p^*\psi_{p'}\,dx
  =\int_0^\infty e^{i(k'-k)x}\,dx
  +\int_0^\infty e^{-i(k'-k)x}\,dx
  =\int_{-\infty}^\infty e^{i(k'-k)x}\,dx,
\]
which, by (15.7), agrees with (21.9).
\SourcePage{88}{91}
According to (5.14), conversion to normalization by an energy delta function
is effected by multiplying $\psi_p$ by
\[
  {\left(\frac{d(p/2\pi\hbar)}{dE}\right)}^{1/2}
  =\frac{1}{\sqrt{2\pi\hbar v}},
\]
where $v$ is the particle's speed at infinity. Thus
\begin{equation}
  \psi_E=\frac{1}{\sqrt{2\pi\hbar v}}\psi_p
  =\frac{1}{\sqrt{2\pi\hbar v}}
  \left(e^{i(kx+\delta)}+e^{-i(kx+\delta)}\right).
  \tag{21.11}
\end{equation}
Notice that each of the two traveling waves into which the standing wave
(21.11) separates has flux density $1/(2\pi\hbar)$.

We may therefore state the following rule for normalizing, to an energy delta
function, the wave function of motion infinite in one direction: write its
asymptotic expression as the sum of two plane waves traveling in opposite
directions, and choose the normalization coefficient so that the flux density
in the wave traveling toward the origin (or away from it) equals
$1/(2\pi\hbar)$.

The analogous argument gives the corresponding rule for wave functions of
motion infinite in both directions. Such a wave function is normalized to an
energy delta function when the sum of the fluxes in the waves traveling
toward the origin from the positive and negative sides of the $x$ axis equals
$1/(2\pi\hbar)$.
